\section{Models and ancillary period consequences}\label{rev:period-models}

\subsection{A simultaneous de Rham calculation}
The elliptic models used in the standard realization are
$E_{s,x}:y^2=4X^3-g_2(x)X-g_3(x)$, with
\[
\begin{array}{c|c|c|c}
s&g_2&g_3&\Delta=g_2^3-27g_3^2\\ \hline
1/6&3&1-2x&108x(1-x)\\
1/4&3(1+3x)&1-9x&729x(1-x)^2\\
1/3&3(1+8x)&1-20x-8x^2&1728x(1-x)^3\\
1/2&3(1-x+x^2)&(1+x)(2-x)(1-2x)/2&729x^2(1-x)^2/4
\end{array}
\]
The discriminants show that the fibers are smooth throughout the
real interval used in the paper except at $x=0$.

There is one connection formula for all four rows. For periods of
$\lambda=dX/y$ and $\mu=X\,dX/y$, it is
\begin{equation}\label{eq:GM}
 \begin{gathered}
 \frac d{dx}\begin{pmatrix}\omega\\h\end{pmatrix}
 =\begin{pmatrix}p&q\\-g_2q/12&-p\end{pmatrix}
     \begin{pmatrix}\omega\\h\end{pmatrix},\\
 \qquad p=-\frac{\Delta'}{12\Delta},\quad
 q=\frac{3(2g_2g_3'-3g_3g_2')}{2\Delta}.
 \end{gathered}
\end{equation}
For completeness, this follows by reduction modulo exact
differentials, without choosing any analytic periods. Put
$P=4X^3-g_2X-g_3$ and $r=-g_2q/12$. The elementary identity
\[
 d\left(\frac{Q}{y}\right)
 =\left(Q_XP-\frac12QP_X\right)\frac{dX}{y^3}
\]
with
\[
 Q_1=-2qX^2+2pX+qg_2/3,\qquad
 Q_2=2pX^2+2rX-pg_2/3-g_2'/12
\]
gives
\[
 \partial_x\lambda=p\lambda+q\mu+d(Q_1/y),\qquad
 \partial_x\mu=r\lambda-p\mu+d(Q_2/y).
\]
The displayed $p,q,r$ are obtained by matching the coefficients of
$X$; their common denominator is the discriminant. Integration over
a flat cycle proves \eqref{eq:GM}.

In the four rows, respectively,
\[
 q=-\frac{\varepsilon_s}{6x(1-x)},\qquad
 \varepsilon_s=(1,1,1,2).
\]
Eliminating $h$ in \eqref{eq:GM} gives
\[
 \omega''-\frac{q'}q\omega'
 -\left(p'+p^2-\frac{q'}q p-\frac{g_2q^2}{12}\right)\omega=0.
\]
Substitution of the factored discriminants reduces this equation to
$x(1-x)\omega''+(1-2x)\omega'-s(1-s)\omega=0$.
This supplies the differential equation used in
Lemma~\ref{lem:period}; the residue at the node determines its
normalization, so no numerical matching of periods is involved.

At $x=1/2$, the invariants are $1728,8000,54000,1728$. The two
less immediate CM assertions follow from the standard quotient by
the point $(0,0)$ of a curve $y^2=x^3+ax^2+bx$:
\[
 (a,b)\longmapsto(-2a,a^2-4b),\qquad
 j=\frac{256(a^2-3b)^3}{b^2(a^2-4b)}.
\]
For $(a,b)=(4,2)$ the quotient has parameters $(-8,8)$ and is
isomorphic over $K$ to the original curve by a scaling with square
$-2$. Composing the 2-isogeny with this isomorphism gives an
endomorphism of degree two, which cannot be an integer
multiplication. The original curve therefore has CM and invariant
$8000$. The curve $y^2=x^3+1$ has CM; translation of its rational
two-torsion point gives parameters $(-3,3)$, whose quotient has
parameters $(6,-3)$ and invariant $54000$. CM is invariant under
isogeny, proving the remaining assertion.

\subsection{What one scalar relation determines}
These consequences describe the value field of a hypothetical
identity; none supplies its missing CM converse.

\begin{proposition}\label{prop:B-value-field}
Write $F=F_s(z)$, $G=\thetaop F_s(z)$ at a standard algebraic
parameter $z\in[-1,1)\setminus\{0\}$. Suppose
$AF+BG=\alpha/\pi$ with $A,B,\alpha\in K$ and
$(A,B)\ne(0,0)$. For $L=cF+dG$ with $Ad-Bc\ne0$, the three
numbers $q_0,\pi,L$ are algebraically independent. Moreover
\[
 \operatorname{trdeg}_K K(\pi,F,G)=2,\qquad
 \ker\bigl(K[P,X,Y]\to K[\pi,F,G]\bigr)
       =\bigl(P(AX+BY)-\alpha\bigr).
\]
\end{proposition}
\begin{proof}
Corollary~\ref{prop:B-relation-space} gives $\alpha\ne0$.
The invertible linear change $(F,G)\mapsto(AF+BG,L)$ and
Theorem~\ref{thm:B-value-independent} show that $q_0,AF+BG,L$
are independent. Replacing the nonzero second coordinate by
$\alpha/\pi$ proves the first claim. The relation also gives
$K(\pi,F,G)=K(F,G)$, of transcendence degree two. Finally
\[
 K[P,X,Y]/\bigl(P(AX+BY)-\alpha\bigr)
       \simeq K[X,Y,(AX+BY)^{-1}],
\]
and evaluation on the algebraically independent pair $(F,G)$ is
injective on this localization.
\end{proof}

\begin{proposition}\label{prop:B-period-field}
If $a\omega^2+b\omega h=\alpha\pi$, with
$a,b,\alpha\in K$ and $\alpha\ne0$, at such a standard parameter,
then $q_0,\omega,h$ are algebraically independent.
\end{proposition}
\begin{proof}
Theorem~\ref{thm:B-value-independent} makes $q_0,R,S$
algebraically independent. The scalar relation gives the invertible
field change
\[
 S=\frac{a\omega+bh}{\alpha},\qquad
 \omega=\frac{\alpha S}{a+bR},\qquad h=R\omega.
\]
Its denominators are nonzero because $\alpha\pi\ne0$.
Thus $K(q_0,\omega,h)=K(q_0,R,S)$ has transcendence degree three.
\end{proof}

\subsection{Functoriality under gauges and pullbacks}
Suppose $\widetilde H(t)=g(t)H(\phi(t))$, with algebraic $g,\phi$,
$\delta=t\,d/dt$, and $D=z\,d/dz$. At an algebraic point $t_0$
where these data are regular, put
$r_\phi=\delta\phi/\phi$ and suppose $g\phi r_\phi\ne0$.
The chain rule is the exact coefficient transport
\begin{equation}\label{eq:general-pullback-transport}
 (A+B\delta)\widetilde H(t_0)
 =g(t_0)\bigl(A+B\delta\log g+B r_\phi D\bigr)H(\phi(t_0)).
\end{equation}
Its triangular coefficient matrix has determinant $r_\phi\ne0$.
It therefore carries algebraic identity lines bijectively, and
preserves all rational projective directions when
$\delta\log g,r_\phi\in\Q$. More invariantly, the condition for
preserving $\mathbf P^1(\Q)$ is that the induced transformation
belong to $\mathrm{PGL}_2(\Q)$. The algebraic factor $g(t_0)$ only
rescales the multiplier.

These qualifications are essential. At a ramification point with
$r_\phi=0$, the nonzero pair $(-\delta\log g,1)$ can have value
zero; for example $g=1$ and $\phi=z_0+(t-1)^2$ give
$\delta\widetilde H(1)=0$ even over a non-CM fiber. Also,
\eqref{eq:general-pullback-transport} concerns specified analytic
branches. Ordinary convergence must be checked for each defining
series, and any independence hypothesis must apply at the new
image parameter $\phi(t_0)$.

\subsection{A general conditional finiteness principle}
\begin{theorem}\label{thm:general-finiteness}
Suppose the data of Section~\ref{sec:general-period-square} descend
to a number field $K_0$, and let $\overline X/K_0$ be a smooth
projective compactification of the base. Let
$t:\overline X\to\mathbf P^1$ be a finite $K_0$-morphism of degree $d$, and
assume $j(E)$ is nonconstant. Fix $e$. If the quadratic-period
independence in Theorem~\ref{thm:general-identity-line} holds for
the non-CM fibers at admissible points with
$[\Q(t(P)):\Q]\le e$, only finitely many such points support a
nonzero rational coefficient identity. At each one the primitive
integer pair is unique up to sign.
\end{theorem}
\begin{proof}
The identity-line theorem forces CM. A fiber of $t$ has degree
$d$, so
\[
 [\Q(j(E_P)):\Q]\le[K_0(P):\Q]
                       \le de[K_0:\Q].
\]
By the ring class field theorem, a singular modulus for the full
endomorphism order of discriminant $\Delta$ has degree
$h(\Delta)$ over $\Q$ \cite{CoxCM}. The finiteness theorem for
imaginary quadratic orders of bounded class number therefore
leaves finitely many possible $j$-invariants. The nonconstant map
$j:\overline X\to\mathbf P^1$ has finite fibers, giving finitely
many points. The uniqueness assertion is unconditional by
Theorem~\ref{thm:general-identity-line}.
\end{proof}
This argument proves finiteness without asserting an effective
algorithm at arbitrary degree. Isotrivial CM families and the
nonadmissible fibers require separate treatment.
